La producción musical ha experimentado en las últimas dos décadas una transformación profunda, impulsada por la democratización de las herramientas digitales de grabación, edición y mezcla. Lo que antes requería estudios profesionales con consolas analógicas y personal técnico especializado, hoy puede realizarse desde un computador personal equipado con un DAW (Digital Audio Workstation) y un conjunto de plugins accesibles a cualquier usuario. Esta accesibilidad ha impulsado el surgimiento de una nueva generación de productores independientes que graban, editan y mezclan su propia música sin haber pasado necesariamente por una formación técnica formal en ingeniería de sonido.

Dentro de este nuevo panorama, la mezcla de audio continúa siendo una de las etapas más determinantes y, a la vez, más exigentes técnicamente dentro del proceso de producción musical. Como señalan De Man, Reiss y Stables en [4], lograr una mezcla equilibrada implica gestionar simultáneamente múltiples decisiones interdependientes de nivel, panoramización, ecualización, compresión y espacialidad, cuya combinación óptima no responde a una fórmula única sino a criterios estéticos, técnicos y artísticos que varían según el género musical y la intención del proyecto. Esta complejidad se acentúa particularmente en la mezcla vocal, dado que la voz humana constituye el elemento central de la mayoría de los géneros populares y su tratamiento exige un equilibrio delicado entre claridad, presencia, calidez y naturalidad, sin que el procesamiento aplicado desvirtúe el carácter propio de la interpretación.

Frente a esta complejidad técnica, la brecha de conocimiento entre productores profesionales y productores independientes o amateurs se ha vuelto cada vez más evidente. Vanka et al. en [1] documentaron, a través de entrevistas, cuestionarios y análisis de foros especializados, que una proporción significativa de productores en contextos de home studio enfrenta la mezcla mediante procesos de ensayo y error, al no contar con una guía técnica clara que oriente sus decisiones. Esta situación no solo prolonga innecesariamente los tiempos de producción, sino que además genera resultados inconsistentes entre proyectos, lo cual limita la calidad final de las obras musicales producidas de manera independiente.

En paralelo a este fenómeno, el campo de la inteligencia artificial aplicada al audio ha experimentado un crecimiento acelerado, ofreciendo nuevas posibilidades para abordar esta brecha técnica. Desde los primeros trabajos de automatización basados en reglas, pasando por el uso de autoencoders profundos para aprender transformaciones de audio directamente desde señales crudas propuesto por Martínez Ramírez y Reiss en [5], hasta arquitecturas más recientes como el sistema de mezcla automática de batería basado en Wave-U-Net desarrollado por Martínez Ramírez, Stoller y Moffat en [2] o el framework generativo MEGAMI presentado por Moliner et al. en [3], la literatura evidencia un interés creciente por automatizar, total o parcialmente, el proceso de mezcla. Sin embargo, la mayoría de estos desarrollos se ha concentrado en la automatización completa de la mezcla instrumental, dejando relativamente poco explorado el caso particular de la mezcla vocal, cuya complejidad tímbrica y variabilidad estilística la hacen significativamente más difícil de modelar que otros instrumentos, tal como lo reconocen los propios Martínez Ramírez y Reiss en [5].

Adicionalmente, trabajos como los de Steinmetz et al. en [7] y Lefford et al. en [8] han señalado que los sistemas completamente automáticos, al operar como cajas negras, tienden a limitar el control creativo del ingeniero y a ignorar el contexto musical, el género y la intención artística que rodean cada decisión de mezcla. En esa misma línea, Venkatesh, Moffat y Miranda en [9] y Mourgela et al. en [10] han explorado, respectivamente, la traducción de descriptores semánticos a configuraciones técnicas y la definición de criterios objetivos de calidad sonora, mientras que Vanka et al. en [11] han avanzado hacia sistemas de transferencia de estilo de mezcla basados en referencias musicales. En conjunto, estos antecedentes sugieren que el camino más prometedor no es la sustitución del criterio humano, sino el desarrollo de sistemas de asistencia inteligente que acompañen y orienten al productor, respetando su libertad creativa.

A partir de este panorama, el presente proyecto se propone explorar el desarrollo de una plataforma de asistencia inteligente orientada específicamente al diseño de cadenas de mezcla vocal, capaz de analizar las características acústicas de una señal de voz y generar propuestas iniciales de procesamiento ajustables por el usuario. Para ello, el presente documento se organiza de la siguiente manera: en primer lugar, se plantea el problema que motiva la investigación y se presenta su justificación; posteriormente, se definen el objetivo general y los objetivos específicos que guiarán el desarrollo del proyecto; a continuación, se expone la pregunta de investigación junto con sus subpreguntas; seguidamente, se desarrolla el marco teórico que sustenta conceptualmente la propuesta, integrando tanto el método de revisión de literatura empleado como los antecedentes más relevantes en el campo; y finalmente, se describe la metodología de desarrollo mediante la cual se pretende construir, entrenar y validar el sistema propuesto.


\section{ANTECEDENTES}
En el enfoque del aprendizaje profundo aplicado a la mezcla, Martínez Ramírez y Reiss, en [5], presentaron en 2017 una de las primeras propuestas formales para explorar redes neuronales profundas como herramienta de mezcla inteligente. En ese trabajo se propuso tratar el procesamiento de stems como una transformación basada en contenido, entrenando autoencoders profundos (DAE) para aprender la cadena de efectos aplicada por ingenieros profesionales a partir de grabaciones crudas. Los experimentos, llevados a cabo con grupos instrumentales de bajo, guitarra, voz y teclados mostraron que el sistema podía aprender aspectos de la transformación espectral, aunque con limitaciones notables en el caso de la voz, ya que su complejidad tímbrica y variabilidad estilística la hace significativamente más difícil de modelar. Este resultado es particularmente relevante para el presente proyecto, ya que pone de manifiesto que la mezcla vocal representa un desafío específico que justifica un sistema especializado.

Siguiendo esa línea, Martínez Ramírez, Stoller y Moffat, en [2], desarrollaron un sistema de mezcla automática de batería basado en una variante de la arquitectura Wave-U-Net, que opera directamente sobre señales de audio en el dominio temporal sin necesidad de representaciones espectrales. El sistema recibe como entrada las pistas individuales de una grabación de batería y produce como salida la mezcla estéreo completa, aprendiendo simultáneamente todas las transformaciones de audio involucradas como ganancia, paneo, ecualización y compresión dinámica sin modelarlas de manera separada. En pruebas de escucha subjetiva con veinte participantes, los resultados del modelo Wave-U-Net fueron perceptualmente indistinguibles de las mezclas realizadas por ingenieros humanos, mientras que el enfoque basado en random forests utilizado como base, obtuvo calificaciones significativamente más bajas. Este estudio resulta una evidencia sólida del potencial del aprendizaje profundo para automatizar la mezcla, pero también resalta que el éxito del modelo depende en gran medida de la consistencia y estructura del material musical utilizado.

Una respuesta a esa limitación fue explorada por Vanka et al. en [1], quienes investigaron de manera sistemática la adopción de herramientas de IA en flujos de trabajo de mezcla musical, prestando especial atención a las expectativas y actitudes de distintos grupos de usuarios. A través de entrevistas con ingenieros profesionales, cuestionarios con semi-profesionales y análisis de foros de internet, los autores confirmaron que las necesidades de los usuarios varían según su nivel de experiencia. Los amateurs buscan soluciones automatizadas que les permitan tratar su música sin necesidad de dominar los aspectos técnicos del proceso de mezcla; los semi-amateurs utilizan las herramientas como punto de partida y como recurso de aprendizaje y los profesionales, que ya cuentan con flujos de trabajo establecidos, prefieren herramientas que asistan y co-creativas que les ofrezcan control y opciones de personalización avanzada. Una observación importante del estudio es que el 54\% de ingenieros entrevistados no utiliza herramientas de IA en sus flujos de trabajo, lo que sugiere que la adopción todavía enfrenta problemas relacionados con la calidad de los resultados, la integración en flujos de trabajo existentes y la confianza en los sistemas automatizados.

Moliner et al. en [3] introdujeron en 2025 MEGAMI, el primer framework generativo para mezcla automática de música multipista. A diferencia de todos los enfoques anteriores que trataban la mezcla como un problema de regresión determinística (prediciendo una única mezcla a partir de un conjunto de pistas), MEGAMI modela la distribución condicional de mezclas profesionales, reconociendo explícitamente que múltiples mezclas válidas pueden existir para las mismas entradas. El sistema opera en un espacio de embeddings de efectos, separando las decisiones de mezcla del contenido musical mediante representaciones aprendidas. En la evaluación objetiva, MEGAMI superó consistentemente a todos los sistemas de línea de base, y en las pruebas de escucha subjetiva se acercó a la calidad de las mezclas realizadas por ingenieros humanos, incluso superándolos en algunos géneros. Esta investigación es de las más recientes en cuanto sistemas automatizados de mezcla, estableciendo un referente técnico importante para el presente proyecto, que busca desarrollar un sistema especializado en la voz con un enfoque asistivo en lugar de completamente automático.

A partir de este mismo problema por evitar sistemas cerrados o poco interpretables, Steinmetz et al. [7] propusieron un enfoque de mezcla multipista basado en una consola de mezcla diferenciable compuesta por efectos neuronales de audio. A diferencia de los modelos que generan directamente una señal final, este sistema estima parámetros de mezcla legibles por el usuario lo que permite revisar y ajustar manualmente el resultado generado. El trabajo resulta importante para el presente proyecto porque plantea una alternativa intermedia entre la automatización completa y la asistencia controlada: el sistema puede aprender convenciones de mezcla a partir de datos reales pero conserva una estructura cercana al flujo de trabajo tradicional del ingeniero basada en parámetros, canales y procesos ajustables. Esta lógica se relaciona directamente con la plataforma propuesta, ya que el objetivo no es entregar una mezcla vocal cerrada sino sugerir una cadena de procesamiento que pueda ser modificada por el usuario.

Desde una perspectiva más conceptual, Lefford et al. [8] introdujeron la noción de sistemas inteligentes de mezcla conscientes del contexto. Los autores señalan que las decisiones de mezcla no dependen únicamente de reglas técnicas, sino también del género musical, la intención emocional, las expectativas del cliente, el entorno de trabajo y la experiencia del ingeniero. Esta investigación permite ampliar la comprensión del problema ya que muestra que una herramienta inteligente no debería limitarse a analizar la señal de audio de manera aislada, sino que también debería considerar las condiciones particulares de uso. Para el presente proyecto este planteamiento es importante porque justifica que la plataforma incorpore información contextual como el género, el tipo de voz, la intención sonora o una referencia musical.

En relación con la comunicación entre usuario y sistema, Venkatesh, Moffat y Miranda [9] estudiaron el uso de word embeddings para la ecualización automática en mezcla de audio. Su investigación propone traducir descriptores semánticos a configuraciones de ecualización, permitiendo que términos como "calida", "brillante", "oscura" o "suave" puedan relacionarse con curvas de EQ. Este antecedente resulta especialmente útil para el desarrollo de una plataforma de asistencia en mezcla vocal, ya que muchos productores y artistas no comunican sus objetivos sonoros mediante parámetros técnicos, sino mediante descripciones perceptuales. Por lo tanto, la incorporación de descriptores semánticos podría facilitar que el usuario exprese la intención de la voz de una manera más intuitiva, mientras el sistema traduce esa intención en una posible recomendación de procesamiento.

Por otra parte, Mourgela et al. [10] aportan una perspectiva basada en el análisis de grandes volúmenes de datos reales de mezclas y masters. Su estudio examina más de 218.000 entradas provenientes de una plataforma web de análisis de audio, considerando variables como loudness integrado, true peak, clipping, compresión, problemas de fase, compatibilidad mono, amplitud estéreo y perfil tonal. Esto es realmente importante para el presente proyecto porque permite justificar la selección de métricas objetivas para evaluar las cadenas generadas por el sistema. Aunque la mezcla vocal mantiene un componente subjetivo, el uso de criterios técnicos medibles permite verificar si las propuestas cumplen condiciones básicas de calidad sonora, como evitar distorsión, controlar la dinámica y mantener un balance tonal coherente.

Finalmente, Vanka et al. [11] desarrollaron Diff-MST, un sistema de transferencia de estilo de mezcla basado en una consola diferenciable, un controlador tipo transformer y una función de pérdida orientada al estilo de producción. El sistema recibe pistas crudas y una canción de referencia para estimar parámetros de efectos dentro de una consola de mezcla, permitiendo generar una mezcla con características similares a la referencia y manteniendo la posibilidad de ajustes posteriores. Este antecedente resulta importante porque introduce la referencia musical como una fuente de información para orientar las decisiones del sistema. En el contexto del presente proyecto, esta idea permite pensar la cadena de mezcla vocal no solo desde las características acústicas de la voz, sino también desde el estilo sonoro que el usuario desea alcanzar.
